\documentclass[12pt, a4paper]{article}

% ─── Packages ────────────────────────────────────────────────────────────────
\usepackage[utf8]{inputenc}
\usepackage[T1]{fontenc}
\usepackage{amsmath}          % Math environments
\usepackage{amssymb}          % Math symbols
\usepackage{graphicx}         % Images
\usepackage{hyperref}         % Clickable links
\usepackage{geometry}         % Page margins
\usepackage{listings}         % Code blocks
\usepackage{xcolor}           % Colors
\usepackage{booktabs}         % Better tables
\usepackage{enumitem}         % Better lists

\geometry{margin=2.5cm}

% ─── Document Info ────────────────────────────────────────────────────────────
\title{\textbf{VimTeX Test Document} \\ \large A Full Feature Tour}
\author{Your Name}
\date{\today}

% ─── Code Style ───────────────────────────────────────────────────────────────
\lstset{
  basicstyle=\ttfamily\small,
  keywordstyle=\color{blue}\bfseries,
  commentstyle=\color{gray},
  stringstyle=\color{orange},
  backgroundcolor=\color{gray!10},
  frame=single,
  breaklines=true,
  numbers=left,
  numberstyle=\tiny\color{gray},
}

% ══════════════════════════════════════════════════════════════════════════════
\begin{document}
% ══════════════════════════════════════════════════════════════════════════════

\maketitle
\tableofcontents
\newpage

% ─────────────────────────────────────────────────────────────────────────────
\section{Introduction}
% ─────────────────────────────────────────────────────────────────────────────

Welcome to this \textbf{VimTeX} test document. This file covers the most
common \LaTeX{} features so you can verify your setup works end-to-end:
compilation, PDF preview, and live reloading.

This document was compiled with \texttt{latexmk} and viewed in
\texttt{Zathura} (or your configured viewer).

% ─────────────────────────────────────────────────────────────────────────────
\section{Math}
\label{sec:math}
% ─────────────────────────────────────────────────────────────────────────────

\subsection{Inline Math}

Einstein's famous equation is $E = mc^2$, and the quadratic formula is
$x = \frac{-b \pm \sqrt{b^2 - 4ac}}{2a}$.

\subsection{Display Math}

The Gaussian integral:
\[
  \int_{-\infty}^{\infty} e^{-x^2} \, dx = \sqrt{\pi}
\]

Maxwell's equations in differential form:
\begin{align}
  \nabla \cdot \mathbf{E}  &= \frac{\rho}{\varepsilon_0} \\
  \nabla \cdot \mathbf{B}  &= 0 \\
  \nabla \times \mathbf{E} &= -\frac{\partial \mathbf{B}}{\partial t} \\
  \nabla \times \mathbf{B} &= \mu_0 \mathbf{J}
                             + \mu_0 \varepsilon_0 \frac{\partial \mathbf{E}}{\partial t}
\end{align}

\subsection{Theorem Environment}

\textbf{Theorem (Pythagorean):} For a right triangle with legs $a$, $b$
and hypotenuse $c$:
\[
  a^2 + b^2 = c^2
\]

% ─────────────────────────────────────────────────────────────────────────────
\section{Text Formatting}
% ─────────────────────────────────────────────────────────────────────────────

\begin{itemize}
  \item \textbf{Bold text} for emphasis
  \item \textit{Italic text} for titles or terms
  \item \underline{Underlined text}
  \item \texttt{Monospace} for code or filenames
  \item \textcolor{blue}{Colored text} using \texttt{xcolor}
\end{itemize}

\begin{enumerate}[label=\arabic*.]
  \item First ordered item
  \item Second ordered item
  \item Third ordered item
\end{enumerate}

% ─────────────────────────────────────────────────────────────────────────────
\section{Tables}
% ─────────────────────────────────────────────────────────────────────────────

\begin{table}[h!]
  \centering
  \caption{A sample comparison table}
  \begin{tabular}{@{} l c c r @{}}
    \toprule
    \textbf{Tool}   & \textbf{Language} & \textbf{Speed} & \textbf{Stars} \\
    \midrule
    Neovim          & C / Lua           & Fast           & 80k+           \\
    VimTeX          & Vim Script        & Fast           & 5k+            \\
    latexmk         & Perl              & Medium         & ---            \\
    Zathura         & C                 & Very Fast      & 3k+            \\
    \bottomrule
  \end{tabular}
\end{table}

% ─────────────────────────────────────────────────────────────────────────────
\section{Code Listing}
% ─────────────────────────────────────────────────────────────────────────────

\begin{lstlisting}[language=Python, caption={A simple Python example}]
def fibonacci(n: int) -> list[int]:
    """Return the first n Fibonacci numbers."""
    seq = [0, 1]
    for i in range(2, n):
        seq.append(seq[-1] + seq[-2])
    return seq[:n]

print(fibonacci(10))
# [0, 1, 1, 2, 3, 5, 8, 13, 21, 34]
\end{lstlisting}

% ─────────────────────────────────────────────────────────────────────────────
\section{References \& Links}
% ─────────────────────────────────────────────────────────────────────────────

You can label sections and refer back to them. For example, see
Section~\ref{sec:math} for the math examples.

\begin{thebibliography}{9}
  \bibitem{vimtex}
    Karl Yngve Lervåg,
    \textit{VimTeX: A modern Vim and Neovim filetype plugin for LaTeX files}.
    \url{https://github.com/lervag/vimtex}

  \bibitem{latexmk}
    John Collins,
    \textit{latexmk: Fully automated LaTeX document generation},
    \url{https://ctan.org/pkg/latexmk}
\end{thebibliography}

% ══════════════════════════════════════════════════════════════════════════════
\end{document}
% ══════════════════════════════════════════════════════════════════════════════
